\chapter{Estado del arte}
\label{chap:estado_arte}

\ph{Este capítulo sitúa el trabajo en su contexto académico y técnico.
Su función es doble: (1) introducir al lector los conceptos teóricos que necesita
para entender el resto del documento y (2) demostrar que el autor conoce
las soluciones existentes y dónde se sitúa la suya.}

\ph{Estructura recomendada: 3--4 secciones temáticas que cubran los pilares
conceptuales del TFG, más una sección final de \emph{Trabajos relacionados}
donde se comparan soluciones existentes con la propia.}

\ph{Regla de oro de citación: toda definición, dato o afirmación que no sea
tuya lleva su cita con \textbackslash cite junto a la propia afirmación, no
acumulada al final del párrafo. Prefiere la fuente primaria (el artículo o
libro que introdujo el concepto) y combínala con publicaciones recientes
(últimos cinco años) y con tipos de fuente diversos (artículos, congresos,
libros, documentación oficial, software). Patrón de ejemplo: «El concepto~X
fue introducido por Apellido et al.~\cite{apellido2024articulo} y desde
entonces se ha aplicado a \dots». Consulta \texttt{bibliography.bib}: incluye
un modelo de cada tipo de entrada.}

\section{Ingeniería de Líneas de Productos Software (SPL)}
\label{sec:ingenieria_spl}

\ph{Primer pilar teórico del trabajo. Por ejemplo: el paradigma, la metodología
o el tipo de sistema sobre el que se construye la solución, apoyado en su
fuente primaria, típicamente un libro de referencia~\cite{apellido2021libro}.}

\ph{Cubre aquí:}
\ph{\begin{itemize}
  \item \ph{Definición y origen del concepto.}
  \item \ph{Por qué es relevante para este TFG.}
  \item \ph{Conceptos auxiliares que el lector necesita conocer.}
\end{itemize}}


\subsection{Líneas de Productos Software}
\label{subsec:spl}

A lo largo de la historia, el desarrollo de software se ha basado tradicionalmente en la construcción de sistemas individuales orientados a resolver necesidades concretas, donde cada nuevo producto software suele desarrollarse como un proyecto independiente, incluso cuando comparte una parte significativa de sus funcionalidades, arquitectura o componentes con otros sistemas existentes. Esta estrategia puede resultar adecuada cuando el número de productos es reducido; sin embargo, hoy en día esto presenta limitaciones importantes para la gran mayoría de organizaciones, ya que surge de forma paralela la necesidad de adaptar un determinado sistema base a una gran cantidad de clientes con problemas ligera o enormemente diferentes.

Debido a este hecho surge la Ingeniería de Líneas de Productos Software, definidas por Clements y Northrop como "\textit{a set of software-intensive systems that share a common, managed set of features satisfying the specific needs of a particular market segment or mission and that are developed from a common set of core assets in a prescribed way}" \cite{clements2002}. Gracias a ellas es posible identificar, gestionar y reutilizar aquellos elementos comunes entre diferentes productos, a la vez que se mantienen explícitamente aquellos aspectos que pueden variar entre las distintas configuraciones posibles, haciendo así factible y mucho más viable la gestión de dichos elementos en lugar de construir cada producto desde cero.

\begin{figure}[H]
    \centering
    \includegraphics[width=0.6\textwidth]{SPL.jpg}
    \caption{Representación de una línea de productos software basada en activos comunes. Extraído de \cite{lineas_producto_software_web}.}
    \label{fig:ejemplo_spl}
\end{figure}

La existencia de elementos comunes y diferencias entre los productos introduce el concepto de variabilidad, definido por van Gurp, Bosch y Svahnberg como «Variability is the ability to change or customize a system» \cite{vanGurp2001}. Esta debe ser identificada, representada y gestionada durante el proceso de desarrollo, ya que permite definir qué elementos pueden formar parte de cada producto y bajo qué condiciones pueden combinarse.

La gestión de esta variabilidad requiere distinguir entre dos actividades complementarias. Por una parte, la ingeniería de dominio (\textit{domain engineering}) se encarga de analizar un determinado dominio, identificar sus elementos comunes y variables, y construir los activos reutilizables que formarán la base de la línea de productos. Por otra parte, la ingeniería de aplicación (\textit{application engineering}) utiliza dichos activos para derivar productos concretos mediante la selección de las características necesarias y la resolución de las variabilidades definidas.

No obstante, este paradigma introduce consigo retos relacionados con la representación y gestión de la complejidad. A medida que aumenta el número de características variables y configuraciones posibles, resulta necesario disponer de mecanismos que permitan gestionar esta variabilidad de forma íntegra y precisa, lo cual da surgimiento a técnicas como los modelos de características.


\subsection{Modelos de características}
\label{subsec:feature_models}

Los modelos de características (\textit{Feature Models}) son considerados uno de los mecanismos más utilizados para representar la variabilidad dentro de las líneas de productos software. Originalmente fueron introducidos como parte del método FODA (\textit{Feature-Oriented Domain Analysis}) por Kang et al., con el objetivo de describir las características comunes y variables existentes dentro de una familia de sistemas software \cite{kang1990}. Benavides et al. los describen de la siguiente forma: "\textit{a feature model represents the information of all possible products of a software product line in terms of features and relationships among them}" \cite{benavides2010}. Estos modelos permiten representar el espacio de configuraciones de una línea de productos mediante las características que pueden formar parte de cada producto y las relaciones existentes entre ellas.

Para entenderlo mejor, estos modelos se pueden representar mediante una estructura jerárquica en forma de árbol, donde cada nodo representa una característica del sistema. Las relaciones existentes entre características permiten establecer dependencias entre ellas, indicando qué elementos son obligatorios, opcionales o alternativas dentro de una configuración concreta. Además, pueden incorporar restricciones adicionales que expresan condiciones que no pueden representarse únicamente mediante la jerarquía del modelo.

\begin{figure}[H]
    \centering
    \includegraphics[width=1\textwidth]{ejemplo_feature_model.png}
    \caption{Ejemplo de modelo de características. Extraído de \cite{benavides2025uvl}.}
    \label{fig:ejemplo_feature_model}
\end{figure}

Aunque los modelos de características tradicionales permiten representar una gran cantidad de escenarios de variabilidad, inicialmente estaban orientados principalmente a expresar relaciones estructurales entre características. Con el aumento de la complejidad de los sistemas software surgió la necesidad de incorporar mecanismos adicionales, como atributos, restricciones aritméticas o diferentes tipos de dominios, que permitan representar escenarios más cercanos a los sistemas reales. A su vez, al incrementar la expresividad de los modelos, también se incrementa la complejidad de los mismos, lo cual hace necesaria la existencia de herramientas capaces de procesar estos modelos automáticamente y verificar propiedades sobre las configuraciones posibles.

\subsection{Análisis automático y motores de razonamiento}
\label{subsec:analisis_automatico_motores_razonamiento}

Una vez definidos los modelos de características como mecanismo para representar la variabilidad de una SPL, surge la necesidad de disponer de técnicas capaces de analizar automáticamente la información contenida en dichos modelos. El análisis automático permite obtener información sobre las configuraciones representadas, comprobar propiedades del modelo y realizar operaciones de razonamiento sin necesidad de explorar manualmente cada una de las posibles variantes.

Los motores de razonamiento (\textit{reasoning engine}) se encargan de realizar este análisis. Procesan los modelos de variabilidad y realizan operaciones como la comprobación de satisfacibilidad, la detección de inconsistencias o la identificación de características que no pueden formar parte de ninguna configuración válida. Por ejemplo, en un restaurante un cliente puede elegir en su primer plato Pasta o Arroz (obligatoriamente uno de ellos), y opcionalmente Salsa. No se puede establecer como restricción que se deba elegir Pasta y Arroz a la vez, ya que no tendría sentido si se ve desde un punto de vista natural. Los motores de razonamiento se encargan de comprobar que esto no ocurra, y lo categorizarían como \underline{no satisfacible}.

\begin{figure}[H]
    \centering
    \includegraphics[width=0.5\textwidth]{ejemplo_fm_propio.png}
    \caption{Ejemplo de modelo de características. Elaboración propia.}
    \label{fig:ejemplo_sat}
\end{figure}

Una de las aproximaciones más utilizadas para realizar análisis automático sobre modelos de características consiste en transformar dichos modelos en fórmulas proposicionales que puedan ser procesadas mediante solucionadores SAT (\textit{Boolean Satisfiability Problem}). El problema SAT consiste en determinar si existe una asignación de valores para un conjunto de variables booleanas que haga verdadera una fórmula lógica dada. En el contexto de las líneas de productos software, las características y restricciones de un modelo pueden representarse mediante variables booleanas y expresiones lógicas equivalentes, permitiendo utilizar solucionadores SAT para responder consultas relacionadas con la validez de configuraciones o la existencia de productos derivados \cite{benavides2010}.

Para aplicar esta técnica, el modelo se traduce normalmente a una fórmula lógica equivalente que conserva las restricciones definidas entre sus elementos. Por ejemplo, una relación obligatoria entre dos características puede expresarse mediante una implicación lógica, mientras que las restricciones de exclusión entre características pueden representarse mediante cláusulas que impidan determinadas combinaciones. Una vez realizada esta transformación, el solucionador SAT puede determinar si existe una configuración válida o identificar aquellas asignaciones que cumplen las restricciones del modelo.

Otra técnica ampliamente utilizada en el análisis automático de modelos de características son los Diagramas de Decisión Binarios (\textit{Binary Decision Diagrams}, BDD). Introducidos originalmente por Bryant como una estructura eficiente para representar funciones booleanas, los BDD permiten representar dichas funciones mediante grafos dirigidos acíclicos cuyos nodos internos corresponden a variables de decisión y cuyos nodos terminales representan los valores posibles de la función \cite{bryant1986}. Así es posible modelar de forma compacta y comprensible el espacio de configuraciones válidas de un modelo de características.

\begin{figure}[H]
    \centering
    \makebox[\textwidth][c]{%
        \includegraphics[width=1.1\textwidth]{sat_vs_bdd.png}
    }
    \caption{Funcionamiento de cada motor de razonamiento. Elaboración propia.}
    \label{fig:sat_vs_bdd}
\end{figure}

La siguiente tabla muestra una comparativa que describe los aspectos principales de ambos \textit{solvers}:

\begin{table}[H]
    \centering
    \begin{tabular}{p{4cm}p{5cm}p{5cm}}
        \toprule
        \textbf{Criterio} & \textbf{SAT} & \textbf{BDD} \\
        \midrule
        Representación del problema &
        Transformación del modelo de características a una fórmula lógica que debe ser evaluada por un solucionador SAT. &
        Construcción de una estructura gráfica que representa las posibles asignaciones de la función booleana. \\
        \midrule
        Tipo de análisis &
        Orientado principalmente a responder consultas concretas sobre la existencia de configuraciones válidas. &
        Permite realizar análisis globales sobre el espacio completo de configuraciones representado. \\
        \midrule
        Resultado obtenido &
        Determina si existe una asignación que satisface las restricciones y, en algunos casos, proporciona una configuración válida. &
        Proporciona una representación compacta del conjunto de configuraciones válidas. \\
        \midrule
        Ventaja principal &
        Gran eficiencia para comprobar propiedades concretas del modelo y resolver problemas de satisfacibilidad. &
        Permite realizar múltiples operaciones de análisis sobre el modelo una vez construido el BDD. \\
        \midrule
        Limitación principal &
        Requiere resolver consultas individuales y no mantiene necesariamente una representación completa del espacio de configuraciones. &
        La construcción del diagrama puede resultar costosa para modelos con elevada complejidad. \\
        \bottomrule
    \end{tabular}
    \caption{Comparativa entre SAT y BDD como técnicas de análisis automático de modelos de características.}
    \label{tab:sat_vs_bdd}
\end{table}


\section{Universal Variability Language (UVL)}
\label{sec:uvl}

A continuación, se ofrecerán unas nociones básicas sobre UVL \cite{uvl}, un lenguaje diseñado para proporcionar una forma común de representar modelos de características y facilitar su procesamiento dentro del ecosistema de ingeniería de líneas de productos software. Entender la estructura y funcionamiento de este lenguaje es esencial, ya que supone la base teórica por la que funciona la solución propuesta. 

El modelo completo utilizado como ejemplo se incluye en el \underline{Apéndice~\ref{app:modelo-uvl}}.


\subsection{Estructura básica de un modelo UVL}
\label{subsec:estructura_basica}

Un archivo UVL está compuesto principalmente por dos bloques:

\begin{itemize}
    \item \texttt{features}: define la estructura jerárquica de las características del modelo, constituyendo el núcleo principal del documento.
    
    \item \texttt{constraints}: define las restricciones lógicas existentes entre características que no pueden expresarse mediante la estructura definida en \texttt{features}. Estas restricciones deben cumplirse para garantizar la consistencia del modelo descrito, ya sea debido a reglas de negocio del producto o a otros requisitos establecidos.
\end{itemize}

Además, UVL permite incorporar dos bloques opcionales en aquellos casos donde sea necesario:

\begin{itemize}
    \item \texttt{include}: permite importar extensiones adicionales, como los \textit{minor levels}.
    
    \item \texttt{imports}: permite importar submodelos de características situados en el mismo directorio, facilitando una organización más modular y limpia de los modelos y sus características.
\end{itemize}

Estos dos últimos bloques deben situarse al principio del archivo UVL.

\subsection{Sintaxis}
\label{subsec:sintaxis}

\subsubsection{Imports}

El mecanismo de importación de UVL permite dividir un modelo de características de gran tamaño en submodelos más pequeños y manejables. Mediante la palabra reservada \texttt{imports}, es posible incorporar otros archivos UVL, facilitando la organización de los modelos y favoreciendo su reutilización. Además, cada submodelo importado puede asociarse a un alias mediante la palabra reservada \texttt{as}, permitiendo acceder a sus características mediante la notación \texttt{alias.FeatureName}.

Los elementos importados pueden utilizarse tanto dentro de la estructura jerárquica de características como en la definición de restricciones, permitiendo mantener los modelos más limpios, modulares y facilitando su mantenimiento y evolución.

Por ejemplo:

\begin{lstlisting}[language={}, caption={Ejemplo de importación de submodelos en UVL.}]
imports
    platform as pl
    security

features
    eShop
        mandatory
            security.Security
        optional
            pl.Platform
\end{lstlisting}

En este caso, los modelos \texttt{platform} y \texttt{security} representan submodelos definidos en archivos UVL independientes situados en el mismo directorio.


\subsubsection{Feature tree}

En UVL, el árbol de características está compuesto por características y grupos de características. Cada modelo debe tener exactamente una característica raíz, a partir de la cual se define la estructura jerárquica del resto de elementos. Cada característica puede contener uno o varios grupos de características asociados.

La relación entre características y grupos se expresa mediante indentación: los grupos se escriben con respecto a su característica padre, mientras que las características se indentan respecto a su grupo correspondiente. Cada característica debe tener un nombre único que permita identificarla dentro del modelo, pudiendo utilizar comillas cuando contiene símbolos especiales o cuando su nombre está formado por varias palabras.

UVL define diferentes tipos de grupos para expresar las relaciones entre características, incluyendo \texttt{mandatory}, \texttt{optional}, \texttt{or} y \texttt{alternative}, pertenecientes al nivel booleano del lenguaje. También incorpora las cardinalidades de grupo (\textit{group cardinality}), que permiten expresar restricciones más avanzadas sobre la selección de características.

Como detalle adicional, cuando una característica lleva mas de una palabra consigo, se suele rodear de comillas dobles \verb|""|.

\subsection{Diseño del lenguaje}
\label{subsec:diseño_lenguaje}

UVL define modelos de variabilidad mediante una estructura jerárquica basada en árboles, donde cada nodo representa una característica y las relaciones entre dichos nodos permiten establecer las reglas de selección entre ellas. El núcleo del lenguaje se basa en diferentes construcciones que permiten especificar cómo pueden o deben seleccionarse las características dentro de un modelo.


\subsubsection{\texttt{mandatory}}

Una característica obligatoria debe ser seleccionada siempre que su característica padre forme parte de una configuración válida. Por ejemplo, en el siguiente modelo, si se selecciona la característica \textit{Sensors}, también debe seleccionarse obligatoriamente \textit{Heart Rate Monitor}.

\begin{lstlisting}[language={}, caption={Ejemplo de característica obligatoria en UVL.}]
Sensors
    mandatory
        "Heart Rate Monitor"
\end{lstlisting}

En este caso, la inclusión de \textit{Sensors} implica necesariamente la selección de \textit{Heart Rate Monitor}.


\subsubsection{\texttt{optional}}

Una característica opcional puede ser seleccionada o no dependiendo de las necesidades de una configuración concreta. Por ejemplo:

\begin{lstlisting}[language={}, caption={Ejemplo de características opcionales en UVL.}]
Sensors
    optional
        GPS
        "Blood Oxygen Sensor" {Accuracy 95}
        Accelerometer
        Gyroscope
\end{lstlisting}

En este caso, si la característica \textit{Sensors} es seleccionada, el usuario puede decidir si incorporar o no características como \textit{GPS}, \textit{Blood Oxygen Sensor}, \textit{Accelerometer} o \textit{Gyroscope}.


\subsubsection{\texttt{alternative}}

Un grupo alternativo indica que, si la característica padre es seleccionada, exactamente una de las características hijas debe formar parte de la configuración. Por ejemplo:

\begin{lstlisting}[language={}, caption={Ejemplo de grupo alternativo en UVL.}]
Display
    alternative
        LCD
        OLED
        "E-Ink"
\end{lstlisting}

En este caso, al seleccionar \textit{Display}, debe escogerse exactamente una de las alternativas disponibles: \textit{LCD}, \textit{OLED} o \textit{E-Ink}.


\subsubsection{\texttt{or}}

Un grupo \texttt{or} permite seleccionar una o más características hijas cuando su característica padre está incluida en una configuración. Por ejemplo:

\begin{lstlisting}[language={}, caption={Ejemplo de grupo OR en UVL.}]
Connectivity
    or
        Bluetooth
        "Wi-Fi"
        LTE
\end{lstlisting}

En este caso, pueden seleccionarse varias opciones simultáneamente, siempre que al menos una de ellas sea seleccionada.


\subsubsection{Restricciones entre características (\textit{Cross-tree Constraints})}

Además de las relaciones jerárquicas definidas mediante la estructura del árbol, UVL permite establecer restricciones lógicas entre características situadas en diferentes ramas del modelo. Estas restricciones, denominadas \textit{cross-tree constraints}, permiten expresar condiciones que no pueden representarse únicamente mediante las relaciones padre-hijo.

Estas restricciones se definen dentro del bloque \texttt{constraints} del archivo UVL y permiten utilizar diferentes operadores lógicos:

\begin{itemize}
    \item Negación mediante \texttt{!A}: expresa que una característica no puede estar seleccionada.
    
    \item Conjunción mediante \texttt{A \& B}: indica que ambas características deben estar seleccionadas.
    
    \item Disyunción mediante \texttt{A | B}: indica que al menos una de las características debe estar seleccionada.
    
    \item Implicación mediante \texttt{A => B}: indica que si una característica A está seleccionada, entonces la característica B también debe estarlo.
    
    \item Equivalencia mediante \texttt{A <=> B}: indica que ambas características deben seleccionarse conjuntamente o ninguna de ellas debe seleccionarse.
\end{itemize}

Por ejemplo, la siguiente restricción establece que si se selecciona la característica \textit{Fall Detection}, también deben seleccionarse \textit{Accelerometer} y \textit{Gyroscope}:

\begin{lstlisting}[language={}, caption={Ejemplo de restricción lógica en UVL.}]
constraints
    "Fall Detection" => (Accelerometer & Gyroscope)
\end{lstlisting}

Otro ejemplo sería:

\begin{lstlisting}[language={}, caption={Ejemplo de restricción de equivalencia en UVL.}]
constraints
    "Voice Assistant" <=> "Wi-Fi"
\end{lstlisting}

En este caso, ambas características deben aparecer juntas en una configuración válida, de forma que la selección de una de ellas implica necesariamente la selección de la otra.

\subsection{Niveles del lenguaje}
\label{subsec:niveles_lenguaje}

En UVL, los niveles del lenguaje permiten separar la expresividad del mismo en diferentes capas. Esta organización permite que usuarios y herramientas puedan decidir qué grado de complejidad desean soportar, además de facilitar la extensión del lenguaje sin incrementar la complejidad para aquellos usuarios que únicamente necesitan modelos sencillos.

Un usuario o herramienta que requiera modelos UVL básicos puede trabajar únicamente con el nivel booleano. Sin embargo, cuando sea necesario representar escenarios más complejos, como atributos numéricos, valores textuales u operaciones sobre ellos, pueden activarse niveles superiores.

UVL define tres niveles jerárquicos de expresividad:

\begin{itemize}
    \item \textbf{Boolean level}: incluye características booleanas y restricciones lógicas. También incluye la incorporación de atributos de diferentes tipos en cada característica.
    
    \item \textbf{Arithmetic level}: incorpora atributos numéricos y operaciones aritméticas.
    
    \item \textbf{Type level}: incorpora características tipadas, que pueden tomar valores enteros, reales o textuales, y permite establecer restricciones sobre dichos valores.
\end{itemize}

Cada uno de estos niveles incluye completamente las capacidades proporcionadas por los niveles anteriores. Además, UVL incorpora niveles menores (\textit{minor levels}), que representan extensiones opcionales que añaden funcionalidades adicionales sin ser necesarias para avanzar entre los diferentes niveles principales.


\subsubsection{Boolean level}

El nivel booleano constituye el nivel principal más básico y representa el núcleo del lenguaje UVL. En él se define la estructura del árbol de características, las relaciones entre ellas (\texttt{mandatory}, \texttt{optional}, \texttt{alternative} y \texttt{or}) y las restricciones lógicas básicas entre características mediante operadores como \texttt{!}, \texttt{\&}, \texttt{|}, \texttt{=>} y \texttt{<=>}.

Por ejemplo:

\begin{lstlisting}[language={}, caption={Ejemplo de modelo UVL utilizando el nivel booleano.}]
include
    Boolean.group-cardinality

features
    Smartwatch
        optional
            GPS
            Display
            "Wi-Fi"
            Bluetooth
            "Water Resistance"
            Accessories
                [0..2]
                    "Leather Strap"
                    "Metal Strap"
                    "Silicone Strap"

constraints
    "Wi-Fi" | Bluetooth
    "Water Resistance" => Display
\end{lstlisting}

En este nivel únicamente se representa la presencia o ausencia de características dentro de una configuración.


\subsubsection{Group Cardinality}

La cardinalidad de grupo (\textit{Group Cardinality}) es un nivel menor perteneciente al nivel booleano. Permite establecer que, dentro de un grupo de características, únicamente pueda seleccionarse un número determinado de elementos comprendido entre un límite inferior y superior.

Por ejemplo, la expresión:

\begin{lstlisting}[language={}, caption={Ejemplo de cardinalidad de grupo en UVL.}]
[0..2]
\end{lstlisting}

indica que, cuando se selecciona la característica Accessories, pueden seleccionarse entre cero y dos de sus características hijas: Leather Strap, Metal Strap y Silicone Strap.


\subsubsection{Arithmetic level}

El nivel aritmético incluye las capacidades del nivel booleano y permite expresar restricciones sobre atributos numéricos asociados a las características. Para ello, se utilizan operadores aritméticos como +, -, * y /, junto con operadores de comparación como ==, !=, <, >, <= y >=.

Además, este nivel incorpora dos niveles menores: \textit{Feature Cardinality} y \textit{Aggregate Functions}.

Por ejemplo:

\begin{lstlisting}[language={}, caption={Ejemplo de atributos y restricciones aritméticas en UVL.}]
include
    Arithmetic.feature-cardinality
    Arithmetic.aggregate-function

features
    Smartwatch
        mandatory
            Connectivity
                or
                    Bluetooth
                    "Wi-Fi"
            Battery
                alternative
                    Standard {Capacity 300}
                    "Long Life" {Capacity 450}
        optional
            "Music Player" {Storage 2}
            "Film Player" {Storage 10}
            "Watch Faces" cardinality [5..20]

constraints
    ("Music Player".Storage + "Film Player".Storage) == 12
    avg(Capacity) >= 300
\end{lstlisting}


\subsubsection{Feature Cardinality}

La cardinalidad de característica (\textit{Feature Cardinality}) es un nivel menor del nivel aritmético que permite indicar un número de instancias que pueden seleccionarse de dicha característica.

Para utilizar esta funcionalidad debe incluirse explícitamente mediante:

\begin{lstlisting}[language={}]
Arithmetic.feature-cardinality
\end{lstlisting}

En el siguiente fragmento con partes omitidas se puede ver la forma en la que se emplea:

\begin{lstlisting}[language={}, caption={Ejemplo de cardinalidad de característica en UVL}]
include
    Arithmetic.feature-cardinality

features
    Smartwatch
      ...
        optional
            "Watch Faces" cardinality [5..20]

constraints
    ...
\end{lstlisting}


\subsubsection{Aggregate Functions}

El nivel menor denominado \textit{Aggregate Functions} permite utilizar funciones agregadas como \texttt{sum()} y \texttt{avg()} dentro del bloque de restricciones.

El siguiente fragmento con partes omitidas muestra dos restricciones agregadas de un modelo cuyas características contienen los atributos numéricos PriceEuros y Calification. La declaración de dichas características se omite para centrar el ejemplo en las restricciones.

\begin{lstlisting}[language={}]
include
    Arithmetic.aggregate-function

features
    ...

constraints
    sum(PriceEuros) > 200
    avg(Calification) > 5
\end{lstlisting}


En este ejemplo, sum(PriceEuros) suma los valores del atributo PriceEuros de las características seleccionadas que lo contienen, y la restricción exige que el resultado sea mayor que 200. Por su parte, avg(Calification) calcula la media de los valores del atributo Calification de las características seleccionadas que lo contienen, y exige que dicha media sea mayor que 5. Estas funciones agregan los valores de los atributos; no cuentan cuántas veces aparecen en el modelo.

\subsubsection{Type level}

El nivel de tipos amplía las capacidades de los niveles booleano y aritmético mediante características tipadas. Además de las características booleanas, permite declarar características de tipo Integer, Real y String, cuyos valores pueden utilizarse en restricciones. Estos tipos se indican delante del nombre de la característica y no deben confundirse con los atributos declarados entre llaves.

Este nivel también incorpora el nivel menor String Constraints, que permite expresar restricciones sobre cadenas mediante funciones como len(), utilizada para obtener su longitud.

Por ejemplo:

\begin{lstlisting}[language={}, caption={Ejemplo de características tipadas y restricciones sobre sus valores en UVL.}]
include
    Type.string-constraints

features
    Smartwatch
        mandatory
            String Color
            Real Weight
            String "Serial Number"

constraints
    len("Serial Number") == 10
    Weight > 30.0
    Color != "Pink"
\end{lstlisting}

En este ejemplo, Color y Serial Number son características de tipo String, mientras que Weight es una característica de tipo Real. La primera restricción exige que el valor de Serial Number tenga exactamente diez caracteres. Las otras dos establecen que el valor de Weight sea mayor que 30.0 y que el valor de Color sea distinto de "Pink".


\section{Ecosistema de herramientas y aplicación en variabilidad software}
\label{sec:ecosistema_herramientas}

A continuación se procederá a hablar sobre herramientas y entornos que permitan trabajar con los modelos de características más allá de su simple representación, las cuales proporcionan capacidades para la creación, gestión y análisis automático de modelos, integrando diferentes técnicas de razonamiento y facilitando su adaptación a nuevos escenarios. Dentro de este ecosistema destaca Flamapy como framework orientado al análisis automático de modelos de características mediante una arquitectura extensible. Asimismo, destaca UVLHub como plataforma destinada a la gestión y compartición de modelos definidos en UVL.

\subsection{Flamapy como framework de análisis de variabilidad}
\label{subsec:flamapy_como_framework}

Flamapy es un framework desarrollado en Python orientado al análisis automático de modelos de características dentro del ámbito de la ingeniería de líneas de productos software \cite{flamapy}. Su objetivo principal es proporcionar un entorno común para trabajar con modelos de variabilidad, facilitando tanto su procesamiento como la aplicación de diferentes técnicas de análisis automático. Para ello, Flamapy puede entenderse desde dos perspectivas complementarias: como una herramienta que permite ejecutar análisis sobre modelos de características de forma directa, y como un framework extensible que proporciona una arquitectura modular para incorporar nuevos modelos, operaciones y técnicas de análisis.

Desde la perspectiva de herramienta, Flamapy proporciona diferentes mecanismos para facilitar su utilización en distintos escenarios. Entre ellos se encuentra una interfaz de línea de comandos (\textit{Command Line Interface}) que permite ejecutar análisis directamente sobre modelos de variabilidad, así como una fachada de Python (\textit{Python Facade}) que simplifica su integración dentro de otros desarrollos software. Además, ofrece interfaces adicionales como servicios REST mediante Swagger y soporte para ejecución en entornos web mediante WebAssembly (WASM), permitiendo utilizar sus capacidades de análisis desde diferentes aplicaciones y contextos.

Como framework, Flamapy sigue una arquitectura modular basada en \textit{plugins}, los cuales permiten extender sus capacidades sin modificar el núcleo principal del sistema. Esta arquitectura hace mucho más sencilla la incorporación de nuevos lenguajes de modelado, metamodelos o técnicas de análisis mediante componentes independientes que pueden integrarse dentro del ecosistema existente. Gracias a este enfoque, es posible adaptar Flamapy a diferentes necesidades relacionadas con la ingeniería de variabilidad, incorporando soporte para nuevos tipos de modelos o herramientas de razonamiento.

Otro de los elementos principales que ofrece Flamapy son las operaciones (\textit{operations}), que representan las diferentes tareas de análisis que pueden ejecutarse sobre modelos de características. Cada operación define una funcionalidad concreta, como la comprobación de satisfacibilidad, la obtención de configuraciones válidas o el cálculo de propiedades del modelo, independientemente del formato en el que este haya sido definido. Esta separación permite reutilizar las mismas operaciones sobre diferentes tipos de modelos siempre que exista una representación compatible dentro del framework.

Por otra parte, las transformaciones (\textit{transformations}) permiten convertir modelos entre diferentes representaciones dentro de Flamapy. Estas transformaciones actúan como mecanismo de conexión entre los distintos metamodelos soportados, permitiendo que los modelos puedan ser procesados por diferentes componentes del framework (o incluso no del framework) independientemente de su representación original. Por ejemplo, el componente \textit{UVLWriter}, que es el que se ha utilizado para descargar los modelos creados por el generador diseñado, permite transformar un modelo de características almacenado internamente en Flamapy a su representación textual en formato UVL. Otras de las transformaciones existentes, por nombrar algunas, son \textit{AFM Reader}, \textit{AFM Writer}, \textit{XML Reader}, \textit{JSON Writer}, \textit{SPLOT Writer} o \textit{FM to BDD}. Todo esto facilita la interoperabilidad entre lenguajes de variabilidad, metamodelos y herramientas de análisis.

\subsection{UVLHub como plataforma de gestión de modelos UVL}
\label{subsec:uvlhub_como_plataforma}

UVLHub es una plataforma orientada a la gestión de modelos de características definidos mediante el lenguaje UVL (\textit{Universal Variability Language}) dentro del ámbito de la ingeniería de líneas de productos software \cite{uvlhub}. Tiene como objetivo proporcionar un repositorio especializado que permita almacenar, organizar, compartir y analizar modelos de variabilidad, facilitando tanto la reutilización de estos artefactos tanto en entornos de investigación como de desarrollo software como la colaboración entre usuarios.

La razón principal de su existencia se debe a la complicada localización de los modelos de características empleados por los usuarios, ya que estos los almacenaban en diferentes repositorios, páginas personales o proyectos independientes, utilizando además diferentes formatos de representación, dificultando su búsqueda, reutilización y análisis. UVLHub sigue los principios de ciencia abierta (\textit{Open Science}), como accesibilidad, transparencia y colaboración, con la finalidad de compartir estos modelos de forma homogénea, haciendo más sencilla la colaboración entre usuarios.

Una de las características más relevantes de UVLHub es su capacidad de integración con otras herramientas y servicios relacionados con la gestión y análisis de modelos de variabilidad. Por una parte, la plataforma se integra con Zenodo, un repositorio abierto de investigación que permite almacenar los conjuntos de datos de forma permanente mediante la generación de identificadores DOI (\textit{Digital Object Identifier}), garantizando que los modelos publicados puedan seguir estando disponibles aunque la plataforma dejase de mantenerse en el futuro. Por otra parte, UVLHub incorpora las capacidades de análisis automático proporcionadas por Flamapy, permitiendo obtener información relevante sobre los modelos almacenados, como el número de características, la validez del modelo o el número de configuraciones posibles que pueden derivarse de ellos.

\begin{figure}[H]
    \centering
    \includegraphics[width=0.8\textwidth]{arquitectura_uvlhub.png}
    \caption{Arquitectura general de UVLHub y sus integraciones principales. Extraído de \cite{wiki_uvlhub}}
    \label{fig:arquitectura_uvlhub}
\end{figure}

Además de sus integraciones externas, UVLHub proporciona diferentes funcionalidades (\textit{features}) en forma de módulos orientadas a cubrir las principales necesidades relacionadas con la gestión de modelos UVL. Algunas de las más relevantes son:

\begin{itemize}
    \item \textbf{dataset}: representa el núcleo principal de UVLHub, permitiendo gestionar el ciclo de vida de los conjuntos de modelos UVL. Esta funcionalidad permite realizar operaciones como la subida de datasets, su almacenamiento junto con los metadatos asociados, la descarga de los modelos y la resolución de identificadores DOI generados mediante la integración con Zenodo.

    \item \textbf{featuremodel}: permite gestionar la información asociada a los modelos de características almacenados dentro de un dataset. Esta \textit{feature} define elementos como el modelo de características, sus metadatos y las métricas asociadas, estableciendo la relación entre los datasets y los archivos UVL almacenados.

    \item \textbf{explore}: proporciona mecanismos para explorar los datasets disponibles dentro de la plataforma mediante diferentes criterios de búsqueda, permitiendo localizar modelos de características publicados previamente.

    \item \textbf{flamapy}: integra las capacidades de análisis automático proporcionadas por Flamapy dentro de UVLHub. Esta integración permite realizar operaciones sobre los modelos UVL, como la validación del modelo o transformaciones hacia otros formatos soportados.

    \item \textbf{zenodo}: permite la comunicación con la API REST de Zenodo para realizar la publicación y almacenamiento permanente de los datasets, proporcionando identificadores DOI asociados a los recursos compartidos.
\end{itemize}

La siguiente imagen representa, a grandes rasgos, los módulos actuales de UVLHub y cómo se complementan entre sí:

\begin{figure}[H]
    \centering
    \includegraphics[width=0.8\textwidth]{modulos_uvlhub.png}
    \caption{Arquitectura general de UVLHub y sus integraciones principales. Extraído de \cite{wiki_uvlhub}}
    \label{fig:modulos_uvlhub}
\end{figure}

Estos módulos son originales de UVLHub. A ellos hay que añadirle también el módulo \textbf{generator}, objetivo principal de este TFG. 

\subsection{Generación automática de modelos de características}
\label{sec:concepto4}

Ya es bien sabido que los modelos de características permiten representar la
variabilidad existente dentro de una línea de productos software. Sin embargo,
a medida que aumenta el número de características, restricciones y posibles
configuraciones, la creación manual de estos modelos puede convertirse en una
tarea costosa y difícil de gestionar.

Por este motivo, la generación automática de modelos de características se ha
utilizado como mecanismo para producir modelos sintéticos destinados, entre
otros usos, a la evaluación y prueba de herramientas de análisis de variabilidad
\cite{segura2011metamorphic,galindo2023llm}.

Estas técnicas permiten obtener conjuntos de modelos con diferentes tamaños,
estructuras y grados de complejidad, facilitando la realización de experimentos
sobre distintos escenarios de evaluación.

\section{Trabajos relacionados}
\label{sec:relacionados}

\ph{Sección clave del estado del arte. Aquí se comparan soluciones, proyectos
o publicaciones que abordan problemas similares al que resuelve el TFG.}

\ph{Empieza con un párrafo que sustente la exhaustividad de la búsqueda:
indica qué bases de datos consultaste (Scopus, Web of Science, Google Scholar,
IEEE Xplore\dots), con qué palabras clave y qué criterios de inclusión y
exclusión aplicaste (p.\,ej.\ publicaciones de los últimos cinco años,
revisadas por pares, con herramienta disponible públicamente). Un párrafo
basta; no hace falta una revisión sistemática formal.}

\ph{Plantilla de redacción:}

\ph{En el ámbito de \dots, destacan los trabajos de \dots, que abordan \dots.
Sin embargo, ninguno cubre \dots, que es precisamente lo que motiva este trabajo.}

\ph{En el ámbito de \dots, los enfoques existentes se centran en \dots.
La diferencia con la propuesta de este TFG es \dots.}

Dentro del ámbito de la ingeniería de líneas de productos software existen diferentes propuestas orientadas a la generación automática de modelos de características. Estas aproximaciones han sido utilizadas principalmente para crear modelos sintéticos destinados a la evaluación de herramientas de análisis de variabilidad, la realización de experimentos controlados o el estudio de la complejidad de diferentes técnicas de razonamiento. Para analizar el estado actual de estas soluciones, se ha tomado como punto de partida el apartado \textit{Related Work} del paper proporcionado por los tutores de este TFG para el desarrollo de un nuevo generador de modelos de características basado en UVL\cite{sundermann2024generating}. En dicho apartado se identifican diferentes propuestas existentes dentro de este ámbito, entre las que se encuentran S.P.L.O.T., FeatureIDE, BeTTy y ETHOM. A partir de esta información, se han realizado búsquedas individuales de cada generador, localizando las páginas web donde se encuentran disponibles sus respectivas publicaciones científicas, como ResearchGate para los trabajos asociados a BeTTy y S.P.L.O.T., ScienceDirect para FeatureIDE y ETHOM, así como las páginas oficiales de los propios proyectos cuando estas proporcionaban información adicional. A partir de estas fuentes se ha extraído la información necesaria para analizar las principales características de cada propuesta, prestando especial atención a sus capacidades de personalización, la expresividad de los modelos generados y su integración dentro de otros ecosistemas de variabilidad software.

\subsection{Generating Feature Models with UVL's Full Expressiveness}

El trabajo desarrollado por Sundermann et al. \cite{sundermann2024generating}
constituye el principal punto de partida del presente TFG.
En dicho trabajo se presenta un generador automático de modelos de características
basado en Universal Variability Language (UVL), capaz de aprovechar los diferentes
niveles de expresividad definidos por el lenguaje.

A diferencia de aproximaciones anteriores centradas principalmente en modelos
booleanos, esta propuesta permite generar modelos que incorporan construcciones
más avanzadas, incluyendo atributos, restricciones aritméticas y elementos
dependientes de tipos. Para ello, los autores definen un sistema configurable que
permite controlar diferentes aspectos relacionados con la estructura y complejidad
de los modelos generados.

El presente TFG toma esta propuesta como base para desarrollar una implementación
integrada dentro del ecosistema Flamapy-UVLHub. Mientras que el trabajo de
Sundermann et al. establece el enfoque de generación y la capacidad de producir
modelos UVL expresivos, este proyecto pretende reflejarlo dentro de una plataforma
existente, incorporando mecanismos de configuración, integración y validación
mediante pruebas automatizadas.

\subsection{S.P.L.O.T.}

S.P.L.O.T. (\textit{Software Product Lines Online Tools}) es una plataforma web orientada al desarrollo y análisis de líneas de productos software, proporcionando diferentes herramientas relacionadas con la gestión de modelos de características \cite{splot}. Entre sus funcionalidades se encuentra un generador automático de modelos de características, diseñado para crear modelos sintéticos que pueden ser utilizados posteriormente en tareas de experimentación y evaluación de técnicas de análisis de variabilidad.

Aunque S.P.L.O.T. supuso una de las primeras aproximaciones ampliamente utilizadas para la generación automática de modelos de características, sus capacidades están principalmente orientadas a construcciones pertenecientes al nivel booleano del lenguaje de variabilidad. Por tanto, permite representar estructuras jerárquicas y restricciones lógicas básicas entre características, pero no contempla elementos más avanzados como atributos numéricos, operaciones aritméticas o restricciones basadas en tipos.

\subsection{FeatureIDE}

FeatureIDE es un framework extensible orientado al desarrollo software basado en características que proporciona soporte para diferentes actividades relacionadas con la ingeniería de líneas de productos software \cite{featureide}. La plataforma permite crear, editar y analizar modelos de características, integrando diferentes herramientas destinadas a facilitar el desarrollo de sistemas configurables.

Además de sus capacidades de gestión y análisis, FeatureIDE incorpora un generador automático de modelos de características que permite crear modelos sintéticos dentro de su propio ecosistema. Sin embargo, al igual que otras propuestas existentes, su generación está principalmente limitada a construcciones pertenecientes al nivel booleano de los lenguajes de variabilidad, por lo que no permite generar modelos que incorporen toda la expresividad proporcionada por lenguajes más recientes como UVL.

\subsection{BeTTy}

BeTTy (\textit{Benchmarking and Testing on the Automated Analysis of Feature Models}) es un framework orientado al benchmarking y evaluación de herramientas destinadas al análisis automático de modelos de características \cite{betty}. A diferencia de aproximaciones más simples, BeTTy permite configurar diferentes parámetros relacionados con la generación de modelos, facilitando la creación de escenarios experimentales adaptados a las necesidades de evaluación.

Una de sus principales aportaciones es la capacidad de generar modelos con diferentes niveles de complejidad estructural y restricciones, permitiendo estudiar el comportamiento de diferentes técnicas de razonamiento sobre escenarios controlados \cite{betty}. Sin embargo, aunque proporciona una mayor flexibilidad respecto a generadores anteriores, sus capacidades siguen centradas principalmente en el nivel booleano del lenguaje de variabilidad. Por ello, permite aumentar la dificultad de análisis de los modelos generados, pero no ampliar su expresividad mediante construcciones como atributos o restricciones aritméticas.

\subsection{ETHOM}

ETHOM (\textit{Evolutionary algoriTHm for Optimized feature Models}) es una
propuesta basada en algoritmos evolutivos orientada a la generación automática
de modelos de características optimizados \cite{ethom}. A diferencia de los
generadores puramente aleatorios, ETHOM formula la generación como un problema
de optimización, utilizando una función de evaluación que permite buscar modelos
con determinadas propiedades.

En particular, ETHOM fue diseñado para generar modelos de características
computacionalmente difíciles de analizar, permitiendo evaluar el comportamiento
de herramientas de análisis de variabilidad bajo escenarios exigentes. Para ello,
el algoritmo recibe diferentes restricciones relacionadas con el modelo a generar
y busca configuraciones que maximicen criterios como el tiempo de ejecución o el
consumo de memoria de las herramientas evaluadas \cite{ethom}.

Esta propuesta supone una evolución respecto a la generación aleatoria de modelos,
ya que permite crear casos de prueba adaptados a objetivos concretos de evaluación.
Sin embargo, al estar orientado a modelos de características tradicionales, no
contempla la generación de construcciones avanzadas propias de lenguajes con
mayor expresividad como UVL.

\subsection{Comparativa de trabajos relacionados}
\label{subsec:comparativa_relacionados}

\ph{Resume la comparación en una tabla: una fila por trabajo (citado con
\textbackslash cite) y 3--4 criterios como columnas, derivados del problema
que resuelve tu TFG (no genéricos). La última fila es tu propuesta: así el
lector ve de un vistazo el hueco que cubre el trabajo.}

\ph{Leyenda: Sí = lo cubre completamente; Parcial = lo cubre con limitaciones;
No = no lo cubre. Comenta la tabla de forma crítica, no descriptiva: justifica
cada «Parcial» y cada «No», identifica el hueco que dejan los trabajos
existentes y cierra con una frase que enuncie la contribución diferencial de
este TFG. Sé honesto: si tu propuesta también tiene algún «Parcial»,
reconocerlo refuerza la credibilidad del análisis. Los criterios (columnas)
elegidos deberán poder respaldarse con evidencia propia en el
Capítulo~\ref{chap:resultados}. El Capítulo~\ref{chap:resultados} retomará
esta tabla para contrastarla con los resultados obtenidos.}

\begin{table}[H]
    \centering
    \renewcommand{\arraystretch}{1.2}
    \begin{tabularx}{\textwidth}{
        >{\raggedright\arraybackslash}X
        >{\centering\arraybackslash}X
        >{\centering\arraybackslash}X
        >{\centering\arraybackslash}X
        >{\centering\arraybackslash}X
    }
        \toprule
        \textbf{Trabajo} &
        \textbf{Generación automática} &
        \textbf{Personalización} &
        \textbf{Expresividad UVL} &
        \textbf{Integración UVLHub} \\
        \midrule
        
        S.P.L.O.T. \cite{splot} &
        Sí &
        Parcial &
        No &
        No \\
        
        FeatureIDE \cite{featureide} &
        Sí &
        Parcial &
        No &
        No \\
        
        BeTTy \cite{betty} &
        Sí &
        Sí &
        No &
        No \\
        
        ETHOM \cite{ethom} &
        Sí &
        Sí &
        No &
        No \\
        
        Sundermann et al. \cite{sundermann2024generating} &
        Sí &
        Sí &
        Sí &
        No \\
        
        \midrule
        
        \textbf{Este TFG} &
        Sí &
        Sí &
        Sí &
        Sí \\
        
        \bottomrule
    \end{tabularx}
    \caption{Comparativa de trabajos relacionados con la generación automática de modelos de características.}
    \label{tab:comparativa_generadores}
\end{table}


% --- Transición al siguiente capítulo ---
\ph{Cierra el capítulo con un párrafo puente de 2--4 líneas: qué queda
establecido aquí y qué pregunta abre el capítulo siguiente. Por ejemplo:
«Identificado el hueco que dejan los trabajos existentes, el
Capítulo~\ref{chap:planificacion} describe cómo se ha organizado el proyecto
para construir la solución que lo cubre: metodología, fases, recursos y
gestión de riesgos».}

El presente capítulo ha establecido los fundamentos teóricos y tecnológicos
necesarios para comprender el contexto del trabajo, así como las principales
soluciones existentes relacionadas con la generación automática de modelos de
características. A partir de este análisis se define el marco sobre el que se
desarrolla la propuesta presentada en este TFG. El siguiente capítulo describe
la planificación seguida durante el proyecto, incluyendo la metodología,
las fases de desarrollo, los recursos empleados y la gestión de riesgos.